\section{La inteligencia artificial se parece más a la manufactura}

Peak dice que la IA se parece más a la manufactura. La afirmación va contra la intuición, pero explica mucho. Muchas personas piensan en la IA como software: el modelo se entrena, el producto sale al mercado y el costo marginal es muy bajo. Sin embargo, los problemas reales que enfrenta una empresa de Agent se parecen más a los de la manufactura: calidad, costo, capacidad de producción, cadena de suministro, entrega, mantenimiento, reproducibilidad y adaptación al escenario de cada cliente.

\begin{figure}[H]
\centering
\includegraphics[width=0.96\textwidth]{ai-like-manufacturing.png}
\caption{La IA se parece más a la manufactura: modelo, capacidad de cómputo, datos, entrega y calidad.}
\end{figure}

\begin{knowledgebox}{Lectura de la figura: la metáfora de la manufactura trata sobre la capacidad de entrega}
En la figura, el modelo es solo la base; la capacidad de cómputo es la capacidad de producción; los datos son la materia prima y el proceso; la entrega es el sitio del cliente, y la calidad es que el resultado sea estable y reproducible. Una empresa de Agent tiene que consolidar todos estos elementos, no solo mostrar la capacidad del modelo.
\end{knowledgebox}

\subsection{Calidad, costo y capacidad de producción}

A la manufactura le importan la tasa de rendimiento, el costo, la capacidad de producción y la cadena de suministro. Con un Agent de IA ocurre algo similar: si una misma tarea se completa de manera estable, si la tasa de fallas puede reducirse, si el costo de inferencia es controlable, si las llamadas a herramientas son confiables, si los permisos y los pagos están integrados y si disminuyen los bloqueos del entorno. Estos problemas no parecen atractivos, pero determinan si el usuario puede depender del producto.

\begin{warningbox}{No trates al Agent solo como una función de software}
Si la tasa de éxito de un Agent es inestable de una ejecución a otra, su costo no es controlable y el entorno lo bloquea con frecuencia, se parece más a una demo artesanal que a un producto entregable. La perspectiva de la manufactura recuerda al equipo que debe atender la estabilidad y la tasa de rendimiento.
\end{warningbox}

\subsection{Ecosistema complementario: pagos, permisos y medidas contra rastreadores}

Peak menciona Cloudflare, la verificación de usuario humano, los sitios web que no permiten el acceso de los Agent y el trabajo de Stripe en pagos para Agent. Esto muestra que el Agent no es solo un problema del modelo: también requiere que todo el ecosistema de internet se adapte en identidad, pagos, permisos, medidas contra rastreadores, auditoría y responsabilidad. Sin estos elementos complementarios, muchas capacidades del Agent se quedan detenidas en el límite del mundo real.

\subsection{Las cuatro carencias de la infraestructura para Agent}

Los pagos, los permisos y las medidas contra rastreadores mencionados antes son solo ejemplos; esta sección organiza esas carencias de forma sistemática. Si se concibe al Agent como un empleado digital capaz de navegar por internet, hacer pedidos, buscar información y operar herramientas en tu nombre, la infraestructura que necesita no se limita a la API del modelo. En primer lugar están la identidad y los permisos: en nombre de quién actúa el Agent, a qué sitios puede acceder y si puede conservar la sesión iniciada. En segundo lugar, los pagos: si el Agent va a comprar algo, contratar un servicio o pagar por usar una herramienta, necesita una interfaz de pago utilizable por máquinas. En tercer lugar, las medidas contra rastreadores y la verificación de usuario humano: muchos sitios web consideran por defecto que el acceso automatizado es un riesgo. En cuarto lugar, la auditoría y la responsabilidad: si el Agent se equivoca, quién responde, cómo se revierte la acción y cómo se rastrea la cadena de operaciones.

\begin{knowledgebox}{Lista de infraestructura para Agent}
\begin{tabularx}{\textwidth}{p{0.22\textwidth}p{0.34\textwidth}X}
\toprule
Infraestructura & Problema que resuelve & Qué pasa si falta \\
\midrule
Identidad/permisos & En nombre de quién actúa el Agent y qué puede hacer & No puede entrar de forma segura a cuentas reales ni a sistemas de negocio. \\
Pagos & Si el Agent puede completar transacciones y compras & Solo puede navegar y sugerir; difícilmente ejecuta el ciclo completo. \\
Colaboración frente a medidas contra rastreadores & Si el sitio web permite el acceso del Agent & Cloudflare, la verificación de usuario humano o el control de riesgos lo bloquean con frecuencia. \\
Auditoría/reversión & Las operaciones se pueden rastrear y los errores se pueden corregir & Las empresas no se atreven a confiar al Agent tareas de alto valor. \\
\bottomrule
\end{tabularx}
\end{knowledgebox}

\begin{importantbox}{La expansión de los Agent no depende solo del modelo}
Cuando el modelo es suficientemente potente, el cuello de botella que decide si los Agent tendrán una expansión masiva se desplaza hacia la identidad, los pagos, los permisos, las medidas contra rastreadores, la auditoría y la responsabilidad: la «infraestructura de internet».
\end{importantbox}

\subsection{Resumen de la sección}

Decir que la IA se parece a la manufactura significa que los productos de Agent, en última instancia, tienen que resolver problemas de calidad, costo, capacidad de producción y entrega. La capacidad del modelo es solo el comienzo; lo decisivo es incorporarse de manera estable a flujos de trabajo reales.
